% ECONOMICS_PROBLEM: {"id": "D82-275", "title": "Mechanism design for nonstandard agents", "jel_code": "D82", "source_id": "OP-0008"}
% !TeX program = xelatex
\documentclass[11pt,a4paper]{article}
\usepackage[margin=23mm]{geometry}
\usepackage{fontspec}
\setmainfont{Latin Modern Roman}
\usepackage{amsmath,amssymb,mathtools}
\usepackage{xcolor,enumitem,booktabs,longtable,array,xurl,hyperref}
\definecolor{muted}{HTML}{53636C}
\hypersetup{colorlinks=true,urlcolor=blue,linkcolor=blue}
\setlength{\parindent}{0pt}
\setlength{\parskip}{5pt}
\newcommand{\R}{\mathbb R}
\newcommand{\E}{\mathbb E}
\newcommand{\N}{\mathbb N}
\newcommand{\ind}{\mathbf 1}
\newcommand{\Var}{\operatorname{Var}}
\newcommand{\Cov}{\operatorname{Cov}}
\newcommand{\supp}{\operatorname{supp}}
\DeclareMathOperator*{\argmin}{arg\,min}
\DeclareMathOperator*{\argmax}{arg\,max}
\newcommand{\entrymeta}[3]{{\sffamily\small\color{muted}\textbf{\texttt{#1}}\quad|\quad #2\par #3\par}}
\newcommand{\Problemref}[1]{\texttt{#1}}
\newcommand{\JELref}[2]{\texttt{#2} (JEL \texttt{#1})}
\begin{document}
\section*{Mechanism design for nonstandard agents}
\textbf{Problem ID:} \texttt{D82-275}\quad \textbf{JEL:} \texttt{D82}\par
% BEGIN ECONOMICS STATEMENT
\label{problem:OP-0008}
\label{atlas:prob:M8}

\noindent\textbf{Original atlas ID:} M8. \textbf{Formulation:} editorial specification of a research family.

\subsection*{Definitions and assumptions}

Let $\tau=(\theta,\alpha)$ be a participant type, where $\theta$ indexes a declared welfare-relevant preference and $\alpha$ indexes cognitive or behavioral characteristics. Types have law $F$ on a specified space. Mechanism $m\in\mathfrak M$ presents a finite decision protocol with outcomes $x\in X$ and payments; context $z$ has known law $P_Z$. A behavioral response model $p\in\mathcal P$ supplies a joint outcome law $p(\mathrm dx\mid m,z,\tau)$, allowing errors and dependence among participants. A welfare evaluator $w_\theta(x)$ and complexity cost $c_\alpha(m,z,x)$ are bounded and fixed before optimization. Feasibility and the chosen participation condition define $\mathfrak M_F$. Neither $w_\theta$ nor complexity cost is automatically identified from inconsistent observed choices.

\subsection*{Formal research question}

Using data from randomized protocols and a specified response-model confidence or identification set $\mathcal P(D)$, characterize mechanisms attaining or approximating
\[
\sup_{m\in\mathfrak M_F}\inf_{p\in\mathcal P(D)}
 \mathbb E_{F,P_Z,p}\!\left[w_\theta(x)-c_\alpha(m,z,x)\right].
\]
Determine which experimental variation identifies responses to mechanisms outside the observed menu, and how much institutional simplification improves the objective. When participants optimally choose cognitive effort, model that effort as a strategic action and require an equilibrium selector; when behavior is an empirical policy, state its invariance restrictions instead. Examine welfare sensitivity to alternative defensible evaluators, not just sensitivity to behavioral parameters, and report outcomes common to all admissible specifications.

\subsection*{Interpretation and scope}

This question requires both positive behavioral prediction and a normative objective. A participant's displayed choice can be an error relative to a stable preference, or it can reveal a context-sensitive preference; those interpretations imply different welfare comparisons. Complexity costs must specify who pays them and how they enter the evaluation. A redesign can change attention, beliefs, or sophistication, so holding fitted behavioral error probabilities fixed across protocols is a substantive assumption. Obvious strategy-proofness is a particular mechanism property with an interpretation for cognitively limited agents; it does not by itself identify welfare or eliminate every misunderstanding. A useful study randomizes protocol complexity and measures understanding, participation, and outcomes alongside choices.

\subsection*{Source audit and status}

Glazer--Rubinstein (1998) is verified as \emph{Motives and Implementation: On the Design of Mechanisms to Elicit Opinions}; its central variation concerns experts' motives, not a general cognitive-error model. Li (2017) and de Clippel (2014) are verified as \emph{Obviously Strategy-Proof Mechanisms} and \emph{Behavioral Implementation}. They supply formal partial answers under their own restrictions. The source atlas overstates the direct cognitive relevance of the first citation and leaves the welfare criterion unspecified. The robust welfare-response program above is editorial, with normative assumptions exposed; it is not an assertion that all behavioral implementation or simple-mechanism design remains unresolved.

\subsection*{References}

\begin{enumerate}

\item[S1] Jacob Glazer and Ariel Rubinstein (1998). \emph{Motives and Implementation: On the Design of Mechanisms to Elicit Opinions}. Journal of Economic Theory 79(2), 157--173. \url{https://www.sciencedirect.com/science/article/abs/pii/S0022053197923851}. \textit{Locator:} abstract; opinion-elicitation model. \textit{Establishes:} Dependence of implementation on experts' motives. \textit{Does not establish:} a general model of cognitive error or complexity cost.

\item[S2] Shengwu Li (2017). \emph{Obviously Strategy-Proof Mechanisms}. American Economic Review 107(11), 3257--3287. \url{https://pubs.aeaweb.org/doi/10.1257/aer.20160425}. \textit{Locator:} abstract; obvious-dominance definition. \textit{Establishes:} A formal strategy-proofness criterion with a behavioral interpretation. \textit{Does not establish:} welfare optimality for every boundedly rational population.

\item[S3] Geoffroy de Clippel (2014). \emph{Behavioral Implementation}. American Economic Review 104(10), 2975--3002. \url{https://www.aeaweb.org/articles?id=10.1257/aer.104.10.2975}. \textit{Locator:} abstract; choice-based implementation framework. \textit{Establishes:} Implementation analysis when context-independent preference maximization is relaxed. \textit{Does not establish:} a unique normative welfare criterion for inconsistent choices.

\end{enumerate}

\subsection*{Common conventions: Source mathematical conventions}

Unless an entry states otherwise, agent and firm sets are finite. State and action spaces are measurable, strategies and outcome maps are measurable, probability laws are specified, and expectations exist for the targets under discussion. Infinite-horizon discount factors lie in $(0,1)$ and discounted objectives are finite. Norms, tolerances and complexity input sizes must be specified for computational comparisons. Constants or primitives that appear locally are inputs, not empirical facts asserted by this catalogue.

\textbf{Identification.} A statistical model $m\in\mathcal M$ implies an observable law $P_m$ and target $\theta(m)$. At law $P$, its identified set is
\[
\Theta(P;\mathcal M)=\{\theta(m):m\in\mathcal M,\ P_m=P\}.
\]
Point identification holds when this set is a singleton, or equivalently when $P_{m_1}=P_{m_2}$ implies $\theta(m_1)=\theta(m_2)$ for all compatible models. Sharp bounds mean bounds attained, or approached when the boundary is not attained, by compatible models. Writing this set is a definition, not a solution: the substantive task is to establish informative restrictions and derive or estimate the set. An unrestricted-confounding model may give only trivial bounds.

\textbf{Inference.} Identification, estimability and valid uncertainty quantification are separate questions. Fix $\alpha\in(0,1)$. A confidence procedure $C_n$ over a declared law class $\mathcal P$ has asymptotic uniform coverage at level $1-\alpha$ when
\[
\liminf_{n\to\infty}\inf_{P\in\mathcal P}\Pr_P\{\theta(P)\in C_n\}\geq1-\alpha.
\]
For set-identified targets the coverage event must be defined separately, for example coverage of the full identified set. If an entry uses identification-indexed models instead of uniquely indexed laws, the same coverage convention is applied over those models. Pointwise limits do not establish uniform validity, and asymptotic validity does not guarantee finite-sample precision. Sample sizes, dependence structures and weak-identification sequences are part of each application.

\textbf{Counterfactuals.} $Y_i(a)$ is a potential outcome under a specified intervention, including its timing, implementation and versions. It is not $Y_i$ conditional on voluntarily choosing $a$. A full intervention vector permits interference; shorthand scalar treatment notation does not silently impose no interference. Assignment, selection, attrition, anticipation and measurement must be specified. A claim about an instrument's local effect is not automatically a population policy effect.

\textbf{Economic equilibrium.} A counterfactual model includes best responses, constraints, feasibility, market clearing, state transitions and expectations consistent with the stated information. When several equilibria exist, either a declared selector is an input or the target is set-valued. Comparisons across policy regimes must use the stated selection convention. Reduced-form relationships are not presumed invariant after an equilibrium-changing intervention.

\textbf{Optimization and welfare.} Optimization is over the stated feasible set. If attainment is not established, $\sup$ or $\inf$ and $\varepsilon$-optimal choices replace an asserted maximizer or minimizer. For example, with value $V=\sup_{a\in\mathcal A}W(a)<\infty$, an $\varepsilon$-optimal set is $\{a\in\mathcal A:W(a)\geq V-\varepsilon\}$ for $\varepsilon>0$. Benefits, costs, risk and health or environmental outcomes can be aggregated only using declared units and conversion weights. Interpersonal utility weights, the treatment of fiscal transfers and foreign profits, discounting, and the behavioral welfare criterion are explicit inputs. A comparative-static derivative additionally requires existence and appropriate differentiability of the selected counterfactual.

\textbf{Sources.} Local labels [S1], [S2], and so forth restart in every question. Locators identify the relevant abstract, model, section or result at the level actually checked; a bibliographic or abstract-level verification is not represented as inspection of an entire proof. Working papers are distinguished from later publications and changing versions. Source summaries are paraphrased. All URL checks and editorial formulations refer to the audit date stated above.
% END ECONOMICS STATEMENT
\end{document}
